\section{Introduction}

When an Indian examination board declares results, lakhs of students open the
result website within the same few minutes. The release time is announced in
advance, almost every request is a look-up of one roll number, and the data
does not change while students read it. It is about as simple as a web
workload can be. Yet result portals keep failing on result day. The chairman
of JEE Advanced 2016 described the portal receiving about \num{10000} requests
per second against a server that could hold \num{2000}~\cite{pagalguy2016jee};
on Maharashtra SSC result day in 2021 schools reported that at least half of
their students could not see results until evening, and the state's education
minister apologised~\cite{ht2021ssc}. The problem is not Indian alone: Kenya's
examination council reported up to \num{300000} hits per second on its results
portal~\cite{standard2025knec}, and in 2025 WAEC in Nigeria withdrew its result
checker over \emph{wrong} results rather than slow ones~\cite{icir2025waec}.

The engineering ingredients that could prevent these failures are well
known: caching dynamic pages and pre-generating them~\cite{challenger1999scalable,labrinidis2000webview},
protecting caches against stampedes~\cite{vattani2015xfetch,nishtala2013memcache},
admission control under overload~\cite{welsh2001seda,elnikety2004gatekeeper},
and the danger of retries that keep an overloaded system down~\cite{bronson2021metastable,huang2022metastable}.
Commercial ``virtual waiting rooms'' are sold for exactly this kind of
event~\cite{cloudflare2026wr,aws2026vwr}. What is missing is a measured,
like-for-like comparison of these options for this workload. The closest
recent studies compare a database-backed and a Redis cache-aside design for a
university portal under fixed concurrency~\cite{putranto2026redis,tofan2026cacheaside};
they do not evaluate pre-computed serving or waiting rooms, do not model a
release-time surge or students who refresh, and do not check that the
answers returned are correct. We found no peer-reviewed evaluation of
user-facing waiting rooms at all.

This paper asks four questions about portals that publish read-only results at
a pre-announced time:
\begin{itemize}
\item \textbf{RQ1 (architecture):} Under the same hardware budget, how far does each common design scale, and at what resource cost?
\item \textbf{RQ2 (waiting rooms):} Does a waiting room improve students' outcomes, and is it fair when students refresh?
\item \textbf{RQ3 (correctness):} Do designs return wrong, missing or stale results under load, after corrections, or when release precedes pre-computation?
\item \textbf{RQ4 (caches and refreshes):} Is cache-aside's cold-start weakness a stampede, and how much do student refreshes amplify load?
\end{itemize}

Our contributions are:
(1)~\textbf{ResultBench}, an open, reproducible testbed with eleven portal
designs under one resource budget, an open-loop load generator with a
parameterised model of how students refresh, and an oracle that checks every
response against ground truth;
(2)~a controlled comparison of database-backed, multi-mirror, cache-aside,
pre-computed key-value and pre-computed static designs;
(3)~to our knowledge, the first controlled evaluation of virtual waiting rooms,
including queue order and ticket handling;
(4)~measurements of correctness failures that only appear on result day;
and (5)~a decomposition of cache-aside misses that shows why the usual
stampede defences do little for per-student keys.
